% Volume IV: Optimization and Computation
% Continuity note for Chapter 20.
% Previous dependencies: Chapters 3-7 provide vector spaces, norms, projections, least squares, quadratic forms, matrix calculus, gradients, Hessians, and Frobenius pairings; Chapters 8-10 provide derivatives, local linearization, Taylor approximation, and local geometry; Chapters 13-18 provide expected losses, likelihoods, cross-entropy, KL divergence, and Fisher geometry.
% New durable concepts: optimization problem statements, objective functions, feasible sets, constraints, minimizers, minimum values, local and global optima, existence conditions, first- and second-order optimality, convex combinations, convex sets, convex functions, epigraphs, Jensen intuition, first-order convexity, subgradients, subdifferentials, proximal intuition, L1 and L2 geometry, regularization, shrinkage, sparsity, and conditioning of objective landscapes.
% Forward hooks: Chapter 21 studies algorithms for solving these problems; Chapter 22 adds constraints and duality; later machine-learning chapters reuse empirical risk, cross-entropy, regularization, and loss landscapes; neural-data chapters reuse convex fitting, sparse decoding, and penalized inverse problems.
% Notation commitments: Decision variables are theta in Theta subset R^d; objectives are L, f, or J; feasible sets are C or Theta; gradients are nabla f(theta); Hessians are nabla^2 f(theta); subgradients are g in partial f(theta); norms use ||.||_2 and ||.||_1.
% Figure plan: no figures are included in this authored source. Proposed raster figures are listed in imagegen-figures/ch20-figure-metadata.md.
\section{Chapter 20: Optimization Geometry and Convexity}
\label{ch:optimization-geometry-convexity}

Learning systems, statistical estimators, and many neural models are built by choosing parameters that make an objective small. Least squares chooses a vector with small residual. Maximum likelihood chooses parameters with small negative log-likelihood. A neural decoder chooses weights with small prediction loss. A dynamical model may choose parameters that make simulated trajectories match recorded activity.

Optimization is the mathematical language for these choices. The central issue is not only how to take a derivative. It is how to state the problem, understand the feasible set, distinguish local from global success, recognize when geometry rules out bad local minima, and see how regularization changes the solution.

\paragraph{Chapter problem.}
How do objective functions, feasible sets, derivatives, convexity, nonsmoothness, and regularization shape the solutions of mathematical problems in neuroscience and machine learning?

\paragraph{Section map.}
Section 20.1 defines optimization problems, minimizers, feasible sets, local and global optima, existence, and first- and second-order conditions. Section 20.2 develops convex sets and convex functions as the main geometry where local information becomes global. Section 20.3 introduces subgradients for nonsmooth convex objectives. Section 20.4 studies regularization as geometry, shrinkage, sparsity, and prior information.

\subsection{20.1 Optimization problems}

\paragraph{Guiding question.}
What exactly is being optimized, over what set, and what does it mean for a solution to be local, global, or well posed?

\paragraph{Beginner-facing intuition.}
An optimization problem is a precise version of ``choose the best allowed object.'' The object might be a scalar threshold, a vector of regression weights, a matrix of synaptic couplings, or millions of neural-network parameters. The word ``allowed'' matters: some parameters may be constrained by nonnegativity, normalization, stability, or physical limits.

The objective assigns a score to each allowed choice. A small score might mean low prediction error, high likelihood after taking a negative log, low energy, or small mismatch with neural data. The feasible set tells us which choices are legal. The solution is not the minimum value itself but an object that attains it.

\paragraph{Formal definitions and notation.}
Let \(\Theta\subseteq\mathbb R^d\) be a feasible set and let
\[
L:\Theta\to\mathbb R
\]
be an objective function. A basic optimization problem is
\[
\boxed{
\min_{\theta\in\Theta} L(\theta).
}
\]
The number
\[
L^\star=\inf_{\theta\in\Theta}L(\theta)
\]
is the infimum or best lower value. A point \(\theta^\star\in\Theta\) is a global minimizer if
\[
L(\theta^\star)\le L(\theta)\qquad \text{for all }\theta\in\Theta.
\]
It is a local minimizer if there exists a radius \(\varepsilon>0\) such that
\[
L(\theta^\star)\le L(\theta)
\]
for all feasible \(\theta\) with \(\|\theta-\theta^\star\|_2<\varepsilon\).

Constraints can be built into \(\Theta\), or written explicitly:
\[
\min_{\theta\in\mathbb R^d}L(\theta)
\quad\text{subject to}\quad
h_j(\theta)=0,\quad g_i(\theta)\le 0.
\]
Chapter 22 develops this constrained form in detail.

\paragraph{Core mechanism: first- and second-order local optimality.}
Assume \(L:\mathbb R^d\to\mathbb R\) is differentiable and \(\theta^\star\) is an unconstrained local minimizer. Then
\[
\boxed{\nabla L(\theta^\star)=0.}
\]
To see why, take any direction \(v\in\mathbb R^d\) and define the scalar function \(\varphi(t)=L(\theta^\star+tv)\). Since \(\theta^\star\) is a local minimizer, \(t=0\) is a local minimizer of \(\varphi\), so \(\varphi'(0)=0\). By the chain rule,
\[
\varphi'(0)=\nabla L(\theta^\star)^\top v.
\]
This must hold for every \(v\), so \(\nabla L(\theta^\star)=0\).

If \(L\) is twice differentiable, the Hessian at an unconstrained local minimizer must be positive semidefinite:
\[
\boxed{v^\top \nabla^2L(\theta^\star)v\ge 0\quad\text{for all }v.}
\]
This follows by applying the one-dimensional second derivative test to \(\varphi(t)\). A point with \(\nabla L=0\) and positive definite Hessian is a strict local minimizer under standard smoothness assumptions. These are local conditions; they do not by themselves guarantee a global minimizer.

Existence also needs assumptions. If \(\Theta\) is compact and \(L\) is continuous, then the Weierstrass theorem guarantees at least one global minimizer. If \(\Theta\) is open or unbounded, or if \(L\) is not continuous, the infimum may fail to be attained.

\paragraph{Concrete worked example.}
Let
\[
L(\theta)=(\theta-3)^2+1,\qquad \Theta=\mathbb R.
\]
Then
\[
L'(\theta)=2(\theta-3),
\]
so the only stationary point is \(\theta^\star=3\). The second derivative is \(L''(\theta)=2>0\), so \(\theta^\star\) is a strict local minimizer. Since the square term is nonnegative,
\[
L(\theta)\ge 1=L(3),
\]
it is also a global minimizer. The minimum value is \(L^\star=1\), while the minimizer is \(\theta^\star=3\).

For a constrained contrast, minimize the same \(L\) over \(\Theta=[0,2]\). The unconstrained minimizer \(3\) is not feasible. The best feasible point is the boundary point \(\theta^\star=2\), even though \(L'(2)=-2\neq 0\). This is why constraints require separate optimality tools.

\paragraph{Computational neuroscience example.}
In a linear neural decoding model, let \(r_i\in\mathbb R^N\) be a population response and \(y_i\in\mathbb R\) a behavioral variable. A decoder chooses \(w\in\mathbb R^N\) to solve
\[
\min_{w\in\mathbb R^N}
\frac{1}{2n}\sum_{i=1}^n (w^\top r_i-y_i)^2.
\]
The decision variable is the weight vector, the objective is mean squared decoding error, and the feasible set is all of \(\mathbb R^N\) unless constraints or regularization are added. A local minimizer satisfies the normal equations from Chapter 5 when the objective is differentiable.

\paragraph{AI and machine-learning example.}
Training a classifier often means minimizing empirical risk:
\[
\min_{\theta\in\mathbb R^d}
\frac{1}{n}\sum_{i=1}^n \ell(f_\theta(x_i),y_i).
\]
Here \(f_\theta\) is the model, \(\ell\) is a loss such as cross-entropy, and \(\theta\) contains the parameters. A deep-network objective may be nonconvex, so a zero gradient can indicate a local minimum, saddle point, or flat region rather than a globally best model.

\paragraph{Common misconceptions and failure modes.}
\begin{itemize}[leftmargin=*]
\item The minimizer \(\theta^\star\) and the minimum value \(L^\star\) are different objects.
\item A zero gradient is necessary for an unconstrained smooth local minimum, but it is not sufficient.
\item Boundary optima need not have zero gradient.
\item A small training objective does not automatically mean good prediction on new data.
\item An optimization problem may be ill posed if no minimizer exists or if many wildly different minimizers fit the same data.
\end{itemize}

\paragraph{Exercises.}
\begin{enumerate}[label=\textbf{20.1.\arabic*.}, leftmargin=*]
\item \textbf{Mechanical.} For \(\min_{\theta\in[0,4]}(\theta-5)^2\), identify the objective, feasible set, minimizer, and minimum value.
\item \textbf{Proof or derivation.} Prove that if \(L:\mathbb R^d\to\mathbb R\) is differentiable and \(\theta^\star\) is an unconstrained local minimizer, then \(\nabla L(\theta^\star)=0\).
\item \textbf{Numerical/computational.} Compute the gradient and Hessian of \(L(\theta_1,\theta_2)=\theta_1^2+4\theta_2^2-2\theta_1\). Find its stationary point and classify it.
\item \textbf{Interpretation.} In a neural decoding problem, what does it mean for the feasible set to include a nonnegativity constraint on decoder weights?
\item \textbf{Misconception/debugging.} A student finds a point with \(\nabla L=0\) and declares the training problem solved. List two additional checks or caveats.
\end{enumerate}

\subsection{20.2 Convex sets and convex functions}

\paragraph{Guiding question.}
What geometric assumptions make every local minimum global and make first-order conditions powerful?

\paragraph{Beginner-facing intuition.}
Convexity means that straight lines stay inside. A set is convex if the line segment between any two feasible points remains feasible. A function is convex if its graph bends upward, so every chord lies above the graph. This rules out isolated bad valleys: if two points are good, everything between them is no worse than the linear interpolation of their objective values.

Convexity is valuable because it turns local information into global information. For differentiable convex functions, the gradient at a point defines a supporting hyperplane below the function. If the gradient is zero, that supporting hyperplane is flat, and no other point can have lower value.

\paragraph{Formal definitions and notation.}
A set \(C\subseteq\mathbb R^d\) is convex if for every \(x,y\in C\) and every \(\lambda\in[0,1]\),
\[
\lambda x+(1-\lambda)y\in C.
\]
The point \(\lambda x+(1-\lambda)y\) is a convex combination of \(x\) and \(y\).

A function \(f:C\to\mathbb R\) on a convex set \(C\) is convex if
\[
\boxed{
f(\lambda x+(1-\lambda)y)
\le
\lambda f(x)+(1-\lambda)f(y)
}
\]
for all \(x,y\in C\) and \(\lambda\in[0,1]\). It is strictly convex if the inequality is strict whenever \(x\neq y\) and \(0<\lambda<1\). It is strongly convex with parameter \(\mu>0\) if
\[
f(y)\ge f(x)+g^\top(y-x)+\frac{\mu}{2}\|y-x\|_2^2
\]
for all \(x,y\) and any valid gradient or subgradient \(g\) at \(x\).

The epigraph of \(f\) is
\[
\operatorname{epi}(f)=\{(x,t):x\in C,\ t\ge f(x)\}.
\]
A function is convex exactly when its epigraph is a convex set.

\paragraph{Core theorem: first-order convex optimality.}
If \(f:C\to\mathbb R\) is differentiable and convex on an open convex set \(C\), then
\[
\boxed{
f(y)\ge f(x)+\nabla f(x)^\top(y-x)
\quad\text{for all }x,y\in C.
}
\]
This is the supporting hyperplane inequality. It says the tangent plane at \(x\) lies below the whole graph.

If \(\nabla f(x^\star)=0\), then for every \(y\in C\),
\[
f(y)\ge f(x^\star)+0^\top(y-x^\star)=f(x^\star),
\]
so \(x^\star\) is a global minimizer. Conversely, if \(x^\star\) is an unconstrained differentiable minimizer in the interior, Section 20.1 gives \(\nabla f(x^\star)=0\). Thus, for smooth convex unconstrained problems, zero gradient and global optimality coincide.

One derivation of the supporting inequality uses the one-dimensional function
\[
\varphi(t)=f(x+t(y-x)),\qquad t\in[0,1].
\]
Convexity of \(f\) implies convexity of \(\varphi\). For a differentiable convex scalar function,
\[
\varphi(1)\ge \varphi(0)+\varphi'(0).
\]
Since \(\varphi'(0)=\nabla f(x)^\top(y-x)\), the inequality follows.

\paragraph{Concrete worked examples.}
The quadratic
\[
f(x)=\frac{1}{2}x^\top Ax-b^\top x
\]
is convex when \(A=A^\top\) is positive semidefinite. If \(A\) is positive definite, it is strictly convex, and the unique minimizer solves
\[
\nabla f(x)=Ax-b=0,
\qquad
x^\star=A^{-1}b.
\]

The scalar function
\[
f(x)=x^4-x^2
\]
is not convex on \(\mathbb R\). Its second derivative is
\[
f''(x)=12x^2-2,
\]
which is negative near \(x=0\). The graph has multiple wells, so local geometry alone does not guarantee global optimality.

\paragraph{Computational neuroscience example.}
Fitting a linear encoding model with squared error and an \(\ell_2\) penalty gives
\[
\min_w \frac{1}{2n}\|Xw-y\|_2^2+\frac{\lambda}{2}\|w\|_2^2.
\]
This objective is convex in \(w\). If \(\lambda>0\), it is strongly convex, so the fitted receptive-field weights are uniquely determined even when the stimulus matrix \(X\) has correlated columns.

\paragraph{AI and machine-learning example.}
Logistic regression with negative log-likelihood is convex in the weight vector:
\[
\sum_{i=1}^n \log(1+\exp(-y_iw^\top x_i)).
\]
Training a multilayer neural network with the same loss is generally nonconvex because the model is nonlinear in its parameters. The loss function can still be optimized successfully, but convex global guarantees no longer apply directly.

\paragraph{Common misconceptions and failure modes.}
\begin{itemize}[leftmargin=*]
\item A convex-looking plot in one slice does not prove the full high-dimensional function is convex.
\item Convexity of the loss in predictions does not imply convexity in parameters of a nonlinear model.
\item Strict convexity gives uniqueness of the minimizer, not necessarily easy numerical conditioning.
\item A nonconvex problem can still be practically solvable; convexity is a guarantee, not the only route to success.
\item Jensen's inequality is about averages inside a convex function; it is not permission to swap nonlinear functions and expectations arbitrarily.
\end{itemize}

\paragraph{Exercises.}
\begin{enumerate}[label=\textbf{20.2.\arabic*.}, leftmargin=*]
\item \textbf{Mechanical.} Decide whether each set is convex: \([0,1]\subset\mathbb R\), the unit disk in \(\mathbb R^2\), and the two-point set \(\{-1,1\}\subset\mathbb R\).
\item \textbf{Proof or derivation.} Prove that the intersection of two convex sets is convex.
\item \textbf{Numerical/computational.} For \(A=\begin{bmatrix}2&0\\0&6\end{bmatrix}\) and \(b=(2,12)^\top\), minimize \(f(x)=\frac12x^\top Ax-b^\top x\). Compute \(x^\star\).
\item \textbf{Interpretation.} Why does adding \(\lambda\|w\|_2^2/2\) with \(\lambda>0\) help make a least-squares problem better posed when stimulus features are highly correlated?
\item \textbf{Misconception/debugging.} A student says, ``Cross-entropy is convex, so every neural network trained with cross-entropy is a convex optimization problem.'' Correct the statement.
\end{enumerate}

\subsection{20.3 Subgradients and nonsmooth optimization}

\paragraph{Guiding question.}
How can optimization use first-order geometry when the objective has corners, kinks, or absolute values?

\paragraph{Beginner-facing intuition.}
Many useful objectives are not differentiable everywhere. The absolute value \(|x|\) has a kink at \(0\). The ReLU \(\max\{0,x\}\) has a kink at \(0\). The hinge loss used in margin classifiers has a kink at the margin. The \(\ell_1\) penalty that encourages sparsity has many kinks.

At a smooth point, the derivative gives the slope of a supporting line for a convex function. At a kink, there may be many supporting slopes. The subgradient collects all slopes of lines that stay below the convex graph.

\paragraph{Formal definitions and notation.}
Let \(f:C\to\mathbb R\) be convex on a convex set \(C\subseteq\mathbb R^d\). A vector \(g\in\mathbb R^d\) is a subgradient of \(f\) at \(x\in C\) if
\[
\boxed{
f(y)\ge f(x)+g^\top(y-x)
\quad\text{for all }y\in C.
}
\]
The set of all subgradients is the subdifferential:
\[
\partial f(x)=\{g:g\text{ is a subgradient of }f\text{ at }x\}.
\]
If \(f\) is differentiable at \(x\), then
\[
\partial f(x)=\{\nabla f(x)\}.
\]

For the absolute value,
\[
\partial |x|=
\begin{cases}
\{1\}, & x>0,\\
[-1,1], & x=0,\\
\{-1\}, & x<0.
\end{cases}
\]
For a vector \(w\in\mathbb R^d\),
\[
\|w\|_1=\sum_{j=1}^d |w_j|,
\]
and the subdifferential is coordinatewise:
\[
g_j=\operatorname{sign}(w_j)\quad\text{if }w_j\neq 0,
\qquad
g_j\in[-1,1]\quad\text{if }w_j=0.
\]

\paragraph{Core mechanism: subgradient optimality and proximal shrinkage.}
For a convex function \(f\), a point \(x^\star\) is a global minimizer if and only if
\[
\boxed{0\in \partial f(x^\star).}
\]
If \(0\in\partial f(x^\star)\), the subgradient inequality gives
\[
f(y)\ge f(x^\star)+0^\top(y-x^\star)=f(x^\star)
\]
for every \(y\). Conversely, if \(x^\star\) minimizes a closed proper convex function under standard conditions, the zero vector is a valid supporting slope at the bottom.

A central worked mechanism is the one-dimensional lasso denoising problem
\[
\min_x \frac12(x-a)^2+\lambda |x|,
\qquad \lambda>0.
\]
The optimality condition is
\[
0\in x-a+\lambda\,\partial |x|.
\]
If \(x>0\), then \(0=x-a+\lambda\), so \(x=a-\lambda\), valid when \(a>\lambda\). If \(x<0\), then \(0=x-a-\lambda\), so \(x=a+\lambda\), valid when \(a<-\lambda\). If \(x=0\), the condition is
\[
0\in -a+\lambda[-1,1],
\]
valid when \(|a|\le\lambda\). Therefore
\[
\boxed{
x^\star=\operatorname{sign}(a)\max\{|a|-\lambda,0\}.
}
\]
This is the soft-thresholding rule. It explains why \(\ell_1\) penalties produce exact zeros.

\paragraph{Concrete worked example.}
Let \(a=3\) and \(\lambda=1\). Then
\[
x^\star=\operatorname{sign}(3)\max\{3-1,0\}=2.
\]
Let \(a=0.6\) and \(\lambda=1\). Then
\[
x^\star=0.
\]
The penalty does not merely shrink small coefficients; it can set them exactly to zero because the subdifferential at the kink contains a whole interval of slopes.

\paragraph{Computational neuroscience example.}
Sparse neural decoding often fits weights by
\[
\min_w \frac{1}{2n}\|Xw-y\|_2^2+\lambda\|w\|_1.
\]
The \(\ell_1\) term encourages the model to use a subset of neurons or features. This can improve interpretability, but it should not be mistaken for proof that only the selected neurons matter biologically. Correlated neural features can cause one of several interchangeable predictors to be selected.

\paragraph{AI and machine-learning example.}
ReLU networks are built from nonsmooth functions. A ReLU is differentiable except at zero, and software libraries choose a conventional subgradient at the kink. Hinge-loss classifiers and \(\ell_1\)-regularized models also rely on subgradient or proximal methods. The ability to optimize nonsmooth objectives is therefore not a corner case; it is routine in machine learning.

\paragraph{Common misconceptions and failure modes.}
\begin{itemize}[leftmargin=*]
\item A nondifferentiable point is not automatically a problem; it may be the optimum.
\item A subgradient is not an arbitrary replacement derivative. It must satisfy the global supporting inequality.
\item Subgradient descent can be slower than gradient descent; stepsize choices matter.
\item The selected support in an \(\ell_1\) model can be unstable under correlated features.
\item Choosing the zero subgradient at a kink in software is a convention for backpropagation, not a statement that the mathematical derivative exists there.
\end{itemize}

\paragraph{Exercises.}
\begin{enumerate}[label=\textbf{20.3.\arabic*.}, leftmargin=*]
\item \textbf{Mechanical.} Compute \(\partial |x|\) at \(x=-2\), \(x=0\), and \(x=5\).
\item \textbf{Proof or derivation.} Prove that if \(0\in\partial f(x^\star)\) for a convex function \(f\), then \(x^\star\) is a global minimizer.
\item \textbf{Numerical/computational.} Apply soft-thresholding with \(\lambda=0.5\) to \(a=(-2, -0.2, 0.4, 3)\) coordinatewise.
\item \textbf{Interpretation.} Why can an \(\ell_1\) penalty make a neural decoder more interpretable but not prove a causal role for selected neurons?
\item \textbf{Misconception/debugging.} A student says, ``ReLU has no derivative at zero, so neural networks with ReLU cannot be trained by gradient methods.'' Explain the role of subgradients and conventions.
\end{enumerate}

\subsection{20.4 Geometry of regularization}

\paragraph{Guiding question.}
How do penalties and constraints change the geometry of solutions, and why do \(\ell_2\) and \(\ell_1\) regularization behave differently?

\paragraph{Beginner-facing intuition.}
Regularization adds information or preference to an optimization problem. It may discourage large weights, encourage sparsity, enforce smoothness, stabilize an inverse problem, or encode a prior belief. Geometrically, it changes which solutions are cheap.

\(\ell_2\) regularization penalizes Euclidean length and tends to shrink weights smoothly. \(\ell_1\) regularization penalizes the sum of absolute values and tends to produce sparse weights. The difference is visible in the shapes of their constraint sets: an \(\ell_2\) ball is round, while an \(\ell_1\) ball has corners on coordinate axes.

\paragraph{Formal definitions and notation.}
A penalized optimization problem has the form
\[
\min_{\theta\in\mathbb R^d}
L(\theta)+\lambda R(\theta),
\qquad \lambda\ge 0,
\]
where \(R\) is a regularizer. A constrained form is
\[
\min_{\theta\in\mathbb R^d}L(\theta)
\quad\text{subject to}\quad
R(\theta)\le c.
\]
Under convexity and suitable regularity conditions, varying \(\lambda\) in the penalized problem corresponds to varying \(c\) in the constrained problem. Chapter 22 makes this relation precise using Lagrange multipliers.

For linear regression with design matrix \(X\in\mathbb R^{n\times d}\) and response \(y\in\mathbb R^n\), ridge regression solves
\[
\min_w \frac12\|Xw-y\|_2^2+\frac{\lambda}{2}\|w\|_2^2.
\]
The first-order condition is
\[
X^\top(Xw-y)+\lambda w=0,
\]
so
\[
\boxed{
w_\lambda=(X^\top X+\lambda I)^{-1}X^\top y
}
\]
when \(\lambda>0\). The added \(\lambda I\) improves conditioning by lifting small eigenvalues of \(X^\top X\).

Lasso regression solves
\[
\min_w \frac12\|Xw-y\|_2^2+\lambda\|w\|_1,
\]
where the nonsmooth \(\ell_1\) geometry can produce exact zeros.

\paragraph{Core mechanism: shrinkage, sparsity, and conditioning.}
Suppose \(X^\top X=Q\Lambda Q^\top\), where \(\Lambda=\operatorname{diag}(\sigma_1^2,\ldots,\sigma_d^2)\). Ridge regression modifies the curvature matrix to
\[
X^\top X+\lambda I
=Q(\Lambda+\lambda I)Q^\top.
\]
The inverse scales directions by
\[
\frac{1}{\sigma_j^2+\lambda}
\]
instead of \(1/\sigma_j^2\). Directions with tiny \(\sigma_j^2\), which are unstable in ordinary least squares, are controlled by \(\lambda\). This is why ridge stabilizes ill-conditioned problems.

\(\ell_1\) regularization has a different mechanism. In the constrained picture, the optimizer finds the lowest loss contour that touches the \(\ell_1\) ball. Because the ball has corners on coordinate axes, tangency often occurs with some coordinates exactly zero. This geometry is the multi-dimensional version of the soft-thresholding rule from Section 20.3.

\paragraph{Concrete worked example.}
For a one-dimensional regression problem with unregularized least-squares objective
\[
L(w)=\frac12(w-4)^2,
\]
ridge regularization solves
\[
\min_w \frac12(w-4)^2+\frac{\lambda}{2}w^2.
\]
The derivative is
\[
(w-4)+\lambda w=0,
\]
so
\[
w_\lambda=\frac{4}{1+\lambda}.
\]
With \(\lambda=1\), the estimate shrinks from \(4\) to \(2\). With \(\lambda=3\), it shrinks to \(1\). Ridge shrinkage is continuous rather than thresholding.

The analogous \(\ell_1\) problem
\[
\min_w \frac12(w-4)^2+\lambda |w|
\]
has solution \(\operatorname{sign}(4)\max\{4-\lambda,0\}\). With \(\lambda=1\), the solution is \(3\). With \(\lambda=5\), the solution is \(0\).

\paragraph{Computational neuroscience example.}
Estimating a receptive field from high-dimensional stimulus features can be ill conditioned because neighboring pixels or time lags are correlated. Ridge regularization stabilizes the estimate by discouraging large coefficients in weakly identified directions. Lasso regularization can select a sparse set of pixels, frequencies, or time lags. Smoothness penalties can be more appropriate when the receptive field is expected to vary gradually across time or space.

\paragraph{AI and machine-learning example.}
Weight decay in neural networks is an \(\ell_2\)-style penalty on parameters. It changes the optimization problem and can improve generalization by discouraging overly large weights, though its exact behavior depends on the optimizer and parameterization. Sparsity penalties, group lasso, and structured regularizers encourage feature selection, channel pruning, or interpretable latent factors. Early stopping can also act like an implicit regularizer by preventing parameters from fully fitting noise.

\paragraph{Common misconceptions and failure modes.}
\begin{itemize}[leftmargin=*]
\item Regularization is not only a trick to prevent overfitting; it also stabilizes ill-posed inverse problems.
\item Larger \(\lambda\) is not always better. It can reduce variance while increasing bias.
\item \(\ell_1\) sparsity is a property of the fitted model, not proof that omitted variables are irrelevant.
\item Penalized and constrained formulations correspond under conditions; they are not identical for every nonconvex problem.
\item Weight decay in adaptive optimizers can differ from classical \(\ell_2\) penalty behavior if implemented differently.
\end{itemize}

\paragraph{Exercises.}
\begin{enumerate}[label=\textbf{20.4.\arabic*.}, leftmargin=*]
\item \textbf{Mechanical.} Write the penalized and constrained forms of ridge regression, naming \(L\), \(R\), \(\lambda\), and \(c\).
\item \textbf{Proof or derivation.} Derive the ridge solution \(w_\lambda=(X^\top X+\lambda I)^{-1}X^\top y\) from the first-order condition.
\item \textbf{Numerical/computational.} For \(L(w)=\frac12(w-6)^2\), compute the ridge solution with \(\lambda=2\) and the lasso solution with \(\lambda=2\).
\item \textbf{Interpretation.} In receptive-field estimation, why might a smoothness penalty be more biologically plausible than a pure sparsity penalty?
\item \textbf{Misconception/debugging.} A student says, ``Regularization only matters when test performance is bad.'' Explain why regularization also matters for identifiability and numerical stability.
\end{enumerate}

\paragraph{Closing bridge.}
This chapter defined optimization problems and the geometry that makes them interpretable. Convexity tells us when local information is globally meaningful, subgradients handle corners, and regularization shapes solutions before algorithms ever enter. Chapter 21 now asks how to compute solutions: gradient descent, momentum, Newton-style curvature methods, and stochastic optimization.

\clearpage
